\section{全局线程池、重排序缓冲区与 pipeline 聚合根}

前几章已经把数值内核（kernel）与 I/O 数据平面（data plane）拆开了；本章只讨论
\texttt{pipeline.cpp}。它不新增任何奇异值分解（Singular Value Decomposition, SVD）公式，
也不重新发明文件格式，而是负责把
\[
    \texttt{testcases}
    \longrightarrow
    \texttt{kernel}
    \longrightarrow
    \texttt{output}
\]
拼成一次端到端执行。换句话说，\texttt{JacobiSvdPipeline} 是应用层聚合根（aggregate root）：
它拥有调度权、关闭协议与错误传播语义，但不偷走 I/O 与 CUDA 内核的领域职责。

当前实现已经不再是“主线程串行计算，写线程异步落盘”的简单模型，
而是演化为显式的生产者-消费者（producer-consumer）并发骨架。
新的骨架是：
\[
    \texttt{input}
    \longrightarrow
    \texttt{dispatch}
    \longrightarrow
    \texttt{global thread pool}
    \longrightarrow
    \texttt{reorder buffer}
    \longrightarrow
    \texttt{single writer}.
\]
这一结构需要同时满足四类要求：较高吞吐、稳定有序的输出、可控的内存驻留，以及不丢失的跨线程错误传播。
这一章就按这四类问题背后的结构来讲，而不再沿着“读到哪段代码就解释哪段代码”的顺序展开。

\subsection{为什么这样设计 pipeline：理论动机、顺序契约与复杂度边界}

先看最朴素的串行方案。若第 $i$ 个 testcase 的主要耗时分解为
\[
    R_i=\text{读取/派发},\qquad
    D_i=\text{解码},\qquad
    C_i=\text{SVD 计算},\qquad
    W_i=\text{写出},
\]
那么串行总耗时近似为
\[
    T_{\mathrm{serial}}
    =
    \sum_{i=0}^{N-1}(R_i+D_i+C_i+W_i).
\]
这种写法实现简单，但它把所有阶段硬串成一条长临界路径（critical path）：磁盘抖动、文本解析、GPU 调度和输出写回彼此阻塞。

本项目更关心的是跨 testcase 的吞吐（throughput），而不是试图通过主机侧多线程人为放大单个矩阵 CUDA kernel 的速度。
因此我们把主线程保留为顺序派发者，把可并行的那部分外包给全局线程池（global thread pool），并让写出阶段异步执行。
于是稳定阶段的节拍更接近
\[
    T_{\mathrm{steady}}
    \approx
    \max(T_{\mathrm{dispatch}},T_{\mathrm{workers}},T_{\mathrm{writer}}),
\]
而不是简单求和。需要强调的是，host 侧线程并发并不自动等价于 device 侧 kernel 的真实并行；
若多个任务共享同一 GPU、同一默认流（default stream）或相同运行时资源，它们仍然可能在设备侧排队。
所以这套 pipeline 优化的是端到端资源组织，而不是单个 SVD 内核的渐近复杂度。

第二个约束是顺序契约（ordering contract）。输入流天然有序：
\[
    X_0,X_1,\ldots,X_{N-1}.
\]
但一旦 testcase 被并发执行，完成顺序就会变成某个排列
\[
    P_{\pi(0)},P_{\pi(1)},\ldots,P_{\pi(N-1)}.
\]
如果直接按完成顺序写文件，那么同一输入在不同运行里可能生成不同顺序的输出；这对实验复现、下游比对和故障定位都很糟糕。
因此系统必须显式恢复
\[
    (U_0,\Sigma_0,V_0),
    (U_1,\Sigma_1,V_1),
    \ldots
\]
这样的单调顺序，而不能把“顺序正确”寄托在线程调度运气上。

第三个约束是内存上界。单个输出数据包（output packet）近似包含
\[
    m_in_i+n_i+n_i^2
\]
个标量，其中 $V_i\in\mathbb{R}^{n_i\times n_i}$ 往往占主导。若没有背压（backpressure），
主线程可以无限提交任务，工作线程可以无限完成结果，最后系统只是在用内存掩盖慢写出。
因此我们把队列窗口参数 \texttt{--queue-capacity} 解释为
\[
    W=\max(\texttt{max\_queued\_results},1),
\]
并让它同时约束 in-flight future 数量与完成结果的堆积规模。若
\[
    E_{\max}=\max_i(m_in_i+n_i+n_i^2),
\]
则长期结果驻留的结构上界可粗略写成
\[
    M_{\mathrm{pending}}=O(W\cdot E_{\max}).
\]
这不是精确的 profiling 结果，而是一个重要的设计结论：资源上界由窗口显式控制，而不是隐含地依赖输入文件长度。

第四个约束是错误诚实性（error honesty）。对于并发系统而言，更危险的失败模式往往不是显式崩溃，
而是部分失败被错误地报告为成功。因此，本章后续的实现细节都围绕同一个目标展开：当输入有序、计算并发、完成乱序而输出又必须有序时，
系统如何在维持吞吐的同时保持确定性（determinism）与失败可审计性（auditability）。

\subsection{全局线程池实现：把并发入口集中起来，而不是到处偷开线程}

\texttt{GlobalThreadPool} 的实现刻意保持最小语义集合。它是一个进程内惰性初始化
（lazy initialization）的全局对象：
\[
    \texttt{global\_thread\_pool()}
    \Rightarrow
    \texttt{static GlobalThreadPool pool(default\_thread\_pool\_size())}.
\]
默认线程数来自 \texttt{std::thread::hardware\_concurrency()}，若系统无法报告，则回退到 $4$。
构造时一次性创建 worker；析构时设置 \texttt{stopping\_}，唤醒所有线程，再逐个 \texttt{join()}。

它的工作线程主循环几乎可以写成三行：
\[
    \text{wait until } \texttt{stopping\_}\lor\neg Q.empty()
    \longrightarrow
    \text{pop one task}
    \longrightarrow
    \text{execute outside lock}.
\]
这段代码的关键特征在于没有引入额外语义。线程池只负责任务排队、唤醒和执行；
它不感知 testcase 编号、矩阵维度和输出顺序，也不持有 Jacobi SVD 的应用语义。
与业务相关的协议均保留在 pipeline 外层。

任务提交通过 \texttt{std::packaged\_task<void()>} 包装，接口是
\[
    \texttt{submit}(f)\longrightarrow \texttt{std::future<void>}.
\]
这样做有两个直接收益。第一，工作线程异常会自动被 future 捕获，主线程可以在合适的位置统一观察；
第二，线程池不需要感知任务结果的数据结构，因为真正的大对象结果并不通过 future 返回，
而是由任务本身提交给输出阶段。future 仅承担任务完成状态的传播职责。

为什么线程池设计为全局对象，而不是每次 \texttt{run()} 都临时构造一组线程？原因主要有三点：

\begin{itemize}
    \item 线程创建与销毁不是零成本；
    \item 并发总量必须有统一入口，否则每一层都觉得自己“只多开一点点线程”，最后总量失控；
    \item pipeline 是应用层协议，不该在每次运行里重复建造同一种基础设施。
\end{itemize}

当然，全局线程池也有明确边界。当前池内任务队列本身没有容量上界，真正的背压发生在 pipeline 外层的 future 窗口。
这意味着“只要所有任务都通过 \texttt{JacobiSvdPipeline::run()} 提交”，资源控制就是可信的；
但若未来其他模块也复用这同一个全局池，就需要重新审视是否要为线程池内部补上 bounded queue、优先级或独立执行器
（executor）语义。

\subsection{流程与重排序缓冲区：主线程派发、工作线程执行、写线程恢复顺序}

从控制流看，\texttt{run()} 的主骨架其实可以压缩成下面这条链：
\[
    \begin{aligned}
        \text{normalize config}
         & \longrightarrow
        \text{dispatch testcase } i
        \longrightarrow
        \text{submit worker future} \\
         & \longrightarrow
        \text{worker computes } P_i
        \longrightarrow
        \text{output.submit}(P_i)
        \longrightarrow
        \text{writer flush}.
    \end{aligned}
\]
这里真正关键的不是阶段数量，而是每个阶段只持有其职责所必需的状态。

\paragraph{第一步：主线程顺序派发。}
主线程保留单游标（single cursor）读取权。这一点十分重要，因为文件输入天然有序、最容易审计，也最不适合在此处引入额外并发复杂度。
对于 \texttt{*.mat} 输入，\texttt{MatDispatchReader} 读取的是
\[
    (\texttt{sequence\_index},\texttt{rows},\texttt{columns},\texttt{buffer})
\]
这样的 \texttt{MatDispatchTask}；其中 \texttt{buffer} 是由 \texttt{cudaMallocHost} 管理的页锁定主机内存
（page-locked host memory, pinned host memory）。主线程只负责把原始网络字节序（network byte order）的 payload
装进这块 pinned buffer，再把任务移动给 worker。对于 \texttt{*.txt} 输入，主线程仍然先解析成
\texttt{io::Matrix}，再把完整矩阵交给线程池，因为文本语法解析天然更适合留在顺序边界上。

\paragraph{第二步：future 窗口充当外层节流阀。}
每次调用 \texttt{pool.submit(...)} 之后，返回的 future 都会进入
\texttt{std::deque<std::future<void>> inflight}。一旦
\[
    |\texttt{inflight}| \ge W,
\]
主线程就立刻对队首 future 做 \texttt{get()}。这一步非常关键：它既避免了无限提交任务，也保证异常不会延迟到输入全部读取完毕后才暴露。
于是 \texttt{--queue-capacity} 不再只是“输出队列容量”，而是整个 pipeline 的 in-flight / reorder window。

\paragraph{第三步：工作线程只做单 testcase 的准备、计算与提交。}
无论输入来自 \texttt{MatDispatchTask} 还是 \texttt{io::Matrix}，worker 端最终都归约为同一工作流：
\[
    \text{prepare Matrix}
    \longrightarrow
    \texttt{KernelStage::execute}
    \longrightarrow
    \texttt{output.submit}(P_i).
\]
在 \texttt{*.mat} 路径里，worker 先执行 \texttt{decode\_dispatch\_task\_matrix}；
随后 \texttt{KernelStage::execute} 验证
\[
    rows>0,\qquad columns>0,\qquad |\texttt{values}|=rows\cdot columns,
\]
再调用 \texttt{one\_sided\_jacobi\_svd}。结果会被重组为
\[
    P_i=(i,s_i,U_i,\Sigma_i,V_i),
\]
其中 $\Sigma_i$ 仍编码成 $1\times n_i$ 的矩阵，而不是裸数组。
这样一来，输出阶段只需要处理一种对象 \texttt{io::Matrix}，从而避免额外的特例分支。

\paragraph{第四步：完成队列吸收并发完成，重排序缓冲区恢复文件顺序。}
\texttt{ReorderBufferOutputStage} 的核心不在于额外引入一个线程，
而在于把输出状态划分为两层结构：
\[
    Q_c=\texttt{completed\_queue\_},\qquad
    B=\texttt{reorder\_buffer\_}.
\]
\texttt{completed\_queue\_} 是工作线程提交完成 packet 的有界 FIFO；\texttt{reorder\_buffer\_} 是写线程私有的
\texttt{unordered\_map<std::size\_t, OutputPacket>}，按 testcase index 暂存乱序结果。工作线程只负责
\[
    P_i\longrightarrow Q_c,
\]
写线程则执行
\[
    Q_c\longrightarrow B
    \longrightarrow
    \text{flush all contiguous } B[\texttt{next\_expected\_index\_}].
\]
只要某个 packet 的 index 正好等于当前期待值，写线程就按固定顺序写出
\[
    U_i,\quad \Sigma_i,\quad V_i,
\]
然后递增 \texttt{next\_expected\_index\_}。于是“完成顺序”和“文件顺序”被显式解耦：前者受调度影响，后者由 index 协议定义。

\paragraph{第五步：背压沿整条链回传。}
\texttt{submit(packet)} 只有在
\[
    |Q_c|<\texttt{queue\_capacity\_}
    \quad\lor\quad
    \texttt{closed\_}
    \quad\lor\quad
    \texttt{worker\_error\_}\ne\varnothing
\]
时才会继续。于是当写线程跟不上时，会形成如下背压链：
\[
    \text{slow writer}
    \rightarrow
    Q_c \text{ full}
    \rightarrow
    \text{worker blocked in submit}
    \rightarrow
    \text{future not ready}
    \rightarrow
    \text{main dispatch blocked}.
\]
这正是期望的行为：慢输出应当通过背压减慢整个系统，而不是以隐式内存膨胀的形式被掩盖。
配合 future 窗口，
重排序缓冲区的长期驻留也被限制在一个有限窗口内，而不会随着输入总长度 $N$ 线性增长。

\subsection{工程健壮性细节：配置归一化、关闭协议、异常传播与不变量}

在并发骨架建立之后，真正决定系统是否具备工程可用性的，是这些看似次要、但实际决定鲁棒性的细节。

\paragraph{其一，运行前先归一化配置，而不是把歧义带进并发区。}
\texttt{run()} 的开头先检查输入、输出路径非空，再通过 \texttt{resolve\_file\_format} 把
\texttt{auto\_detect} 落成具体格式，只接受 \texttt{.mat} 和 \texttt{.txt}。若用户启用
\texttt{layout\_transpose\_auto\_tune} 且布局模式仍是 \texttt{auto\_select}，pipeline 会先做一次
阈值自动调优（auto-tuning），把推荐阈值写回本次运行的 kernel 配置，然后关闭本次执行中的调优标志。
这一步必须发生在并发启动之前；否则多个 worker 各自执行 benchmark，将导致同一运行内的语义混乱和资源争用。

\paragraph{其二，关闭协议必须区分成功路径与失败路径。}
成功路径上，主线程读完整个输入、观察完所有 future 后调用
\[
    \texttt{output.close\_success(submitted\_count)}.
\]
写线程排空完成队列与重排序缓冲区后退出；随后系统检查
\[
    \texttt{written\_count\_}=\texttt{expected\_count\_}.
\]
若不相等，则说明存在缺失 packet，例如某个 testcase 永远没有抵达，使重排序缓冲区无法越过空洞。
此时应直接抛出错误，而不是将部分写出的文件误判为成功。失败路径上，任一 worker 异常或主线程读入失败都会触发
\[
    \texttt{output.close\_error(exception)}.
\]
这会设置 \texttt{worker\_error\_}、\texttt{closed\_} 与 \texttt{abort\_}，唤醒所有等待者并快速收尾。
当前实现选择的是 fail-fast 语义，而不是“跳过异常 testcase 后继续写出后续结果”。

\paragraph{其三，异常传播有两条通道，但最终只允许一个结果：\texttt{run()} 不能错误报告成功。}
工作任务内部的解码失败、矩阵验证失败、CUDA runtime 错误、\texttt{output.submit} 错误，
都会进入对应 future，由主线程在 \texttt{future.get()} 时接住并保存为第一个 worker 错误。
写线程内部的目录创建失败、writer 构造失败、重复 index、写文件失败等异常，则保存到
\texttt{worker\_error\_}，再由后续的 \texttt{submit()} 或关闭路径重抛。于是错误传播图可以写成：
\[
    \begin{aligned}
        \text{worker failure}
         & \rightarrow
        \texttt{future}
        \rightarrow
        \texttt{run fails}, \\
        \text{writer failure}
         & \rightarrow
        \texttt{worker\_error\_}
        \rightarrow
        \texttt{submit/close rethrows}
        \rightarrow
        \texttt{run fails}.
    \end{aligned}
\]
这两条通道共同保证：一旦任意线程认为结果不再可信，主线程最终就不能报告成功。

\paragraph{其四，资源释放依赖 RAII（Resource Acquisition Is Initialization），但错误报告依赖显式协议。}
线程池析构负责停止并 \texttt{join()} worker；输出阶段析构调用 \texttt{close\_noexcept()} 作为兜底，
避免 C++ 在栈展开（stack unwinding）期间再次抛异常。也就是说，析构函数负责“别泄漏资源”，
而真正负责“把错误诚实地告诉调用者”的，仍然是显式的 \texttt{close\_success()} 与 \texttt{close\_error()}。
这对应一条经典但可靠的工程原则：析构函数负责资源回收，而核心语义应通过显式协议表达。

\paragraph{其五，不变量（invariants）比技巧更重要。}
只要没有异常，本章设计依赖的核心不变量其实很少：

\begin{itemize}
    \item 每个 testcase 都获得唯一且单调递增的 \texttt{testcase\_index}；
    \item 每个成功 worker 恰好产生一个 \texttt{OutputPacket}；
    \item packet 的 index 在输出阶段不会被修改；
    \item 完成队列只影响抵达顺序，不定义文件顺序；
    \item 只有当 \texttt{next\_expected\_index\_} 命中时，写线程才会真正落盘；
    \item 每个 testcase 固定写出三张矩阵：$U/\Sigma/V$；
    \item 任意线程异常最终都能沿 future 或 \texttt{worker\_error\_} 回到主线程。
\end{itemize}

并发系统最容易失控的情形，是状态过多、例外过多且协议在不同层之间不断漂移。
当前实现之所以仍然可分析，正是因为它把关键状态压缩成了：
序号、future 窗口、完成队列、重排序缓冲区、下一个期望 index 和异常指针。
只有在状态空间足够受限时，正确性才可能被系统性推理。

\subsection{当前实现的边界与可演化方向}

在上述设计下，当前实现已经为顺序输入、并发执行与有序输出建立了较稳定的协议边界；
但它仍然保留了若干明确的系统边界。这些边界不应被简单理解为缺陷，
而应被视为后续演化时必须显式处理的设计决策点。

\paragraph{第一，线程池大小尚未成为 pipeline 配置的一部分。}
当前全局线程池的 worker 数由硬件并发数决定，而不是由 CLI 或 \texttt{PipelineConfig} 显式控制。
这种设计降低了配置表面积，但也意味着用户当前只能控制 in-flight 窗口，
而不能直接限制主机侧工作线程数。若未来 profiling 表明线程数本身成为关键变量，
则可以考虑将其提升为进程级配置；不过这会进一步牵涉全局对象初始化与生命周期语义。

\paragraph{第二，文本输入仍然停留在主线程解析。}
\texttt{*.mat} 路径已经具备较适合派发的任务形态，而 \texttt{*.txt} 路径仍然在主线程完成解析，
随后再把 \texttt{io::Matrix} 提交给线程池。当前做法强调错误定位与控制流清晰，
但若文本输入在某些工作负载下成为瓶颈，则可能需要进一步考虑并行文本解析。
这一演化会同时带来更复杂的输入游标协议、错误定位与行列一致性诊断问题。

\paragraph{第三，pipeline 级并发不等价于显式 CUDA stream 流水线。}
当前 worker 任务虽然可以并发提交主机侧工作，但这并不等价于显式的
\[
    \text{host decode}
    \parallel
    \text{host-device copy}
    \parallel
    \text{kernel}
    \parallel
    \text{device-host copy}
\]
多阶段 CUDA stream 流水线。若要进一步挖掘设备侧重叠，
则需要 kernel 层暴露异步接口、stream 与事件（event）语义，并由更外层协议显式管理生命周期；
这不应在当前 pipeline 层中隐式引入。

\paragraph{第四，报告接口仍偏向语义报告，而非完整遥测。}
\texttt{PipelineReport} 当前主要报告 testcase 数量、输出矩阵数量、sweep 总数以及布局转置的运行时配置。
它尚未覆盖每阶段耗时、future 阻塞次数、完成队列等待时间、重排序缓冲区峰值占用等性能指标。
如果未来需要系统地分析吞吐瓶颈，更合理的方向是新增独立 profiling/telemetry 结构，
而不是让现有报告对象持续承担调试与性能观测职责。

\paragraph{第五，边界条件的当前策略是保守的 fail-fast 语义。}
当前实现对缺失 packet、写线程失败、worker 异常和关闭阶段异常都采取快速失败策略，
而不尝试在输出文件中保留部分成功结果，也不定义跳过坏 testcase 后继续执行的语义。
这一选择有利于保持实验结果的可审计性，但也意味着若未来要支持部分恢复（partial recovery）、
按 testcase 容错或断点续跑，则必须同步重写输出格式、关闭协议和报告语义。

\subsection{小结}

本章将 \texttt{pipeline.cpp} 重新定位为应用层聚合根（aggregate root）：它既不实现底层 I/O，
也不参与单矩阵 SVD 的数值推导，而是负责把顺序输入、并发执行与有序输出组合成一个可验证的执行协议。
围绕这一定位，系统采用全局线程池来集中管理主机侧并发，
采用 future 窗口来限制 in-flight 任务数量，
采用完成队列与重排序缓冲区来恢复输出顺序，
并通过显式背压、关闭协议与异常传播保证资源边界与失败语义的一致性。

因此，本章的重点并不在于“引入了线程池和重排序缓冲区”这一实现事实本身，
而在于这些机制如何共同定义了一个有限、确定且可审计的应用层边界。
对于本项目而言，\texttt{pipeline.cpp} 的价值并不是增加并发结构的数量，
而是在输入有序、计算并发、完成乱序而输出仍需稳定有序的约束下，
给出一个在吞吐、正确性与工程鲁棒性之间更平衡的系统组织方式。
